\documentclass[11pt]{article}

\usepackage{enumitem}
\usepackage{needspace}
\usepackage{amsmath,amssymb}
\usepackage{times}
\usepackage[normalem]{ulem}
\usepackage{xcolor}
\usepackage[T1]{fontenc}
\usepackage{pdfpages}
\usepackage{hyperref}
\definecolor{responseblue}{RGB}{0,0,180}
\definecolor{revisionred}{RGB}{220,0,0}
\hypersetup{colorlinks=true,urlcolor=responseblue,linkcolor=responseblue}
\setlength{\emergencystretch}{4em}
\Urlmuskip=0mu plus 1mu\relax

\newcommand{\reviewerheading}[1]{%
  \par\vspace{6mm}\noindent\textbf{\uline{#1}}\par\vspace{3mm}}
\newcommand{\response}[1]{%
  \Needspace{3\baselineskip}%
  \par\noindent{\color{responseblue}\textbf{Response to #1:}}\par
  \vspace{1mm}\color{responseblue}}
\newcommand{\finishresponse}{\par\color{black}\vspace{5mm}}
\newcommand{\revloc}[1]{%
  \textcolor{revisionred}{\textbf{\uline{#1}}}}

\textwidth 6.4in
\textheight 8.5in
\topmargin -0.3in
\evensidemargin 0in
\oddsidemargin 0in
\linespread{1.1}

\begin{document}

\begin{center}
  {\bfseries Responses to the Comments of the Editor}\\[5pt]
  {\itshape \textbf{Title:} Chain-Structured Binary Neural Architecture Search for Efficient Visual Recognition across Platforms}\\
  {\itshape \textbf{Authors:} Rui Wan, Dong Wang, Yiwen Liang, and Mingrun Wei}\\
  {\itshape \textbf{Journal:} Image and Vision Computing}\\
  {\bfseries Manuscript Number (Full):} IMAVIS-D-26-01530R1\\
  {\bfseries Manuscript Number (Abbreviated):} IMAVIS-D-26-01530
\end{center}

\vspace{3mm}
We sincerely thank the editor for the careful assessment of our manuscript and the constructive comments. These comments have helped us identify areas where the methodology, experimental evaluation, and presentation require clarification and improvement. We have carefully revised the manuscript in response and provide point-by-point responses below.

The editor's comments are presented in \textbf{bold black text}. Each comment is followed by a space for its response. Responses will be set in \textcolor{responseblue}{\textbf{blue}}, and revision locations will be marked in \textcolor{revisionred}{\textbf{\uline{red and underlined text}}}.

To allow the editor and reviewers to consult the supplementary material without downloading a separate file, a complete copy is appended after the references in this response letter. The appended supplementary material retains its own layout, starts at page 1, and has independently numbered sections, figures, tables, equations, and references.

\reviewerheading{Comments of the Editor.}

\begin{enumerate}[leftmargin=0pt,label={},itemsep=5mm]
\bfseries

\item[]\textbf{Major Comment 1:} \textbf{Inconsistent and incomplete accounting of the headline search-cost claim.}
The headline claim of ``0.48 GPU-hours per platform'' is not supported by consistent accounting across the manuscript.
\begin{enumerate}[label=(\alph*),leftmargin=2em,itemsep=2mm]
  \item Training CIF-Net requires accuracy labels for 750 architectures, each trained for 25 epochs, at a stated cost of 1363.2 GPU-hours (Sec. 4.2.1). This is about 2.8 times the 480 GPU-hours reported for the entire NAS-BNN search, which the manuscript describes as prohibitively expensive (Secs. 1 and 2.1). The CIF-Net labeling cost is excluded from the search-time columns of Tables 3 and 4, whereas the NAS-BNN cost is presented as a fully inclusive number. Although the amortization argument in Sec. 5.5 is reasonable, the abstract and Table 3 currently compare an online-only cost against a fully inclusive baseline cost. Under complete accounting, FlashNAS-BNN costs more than NAS-BNN for fewer than approximately three target platforms. The manuscript should report the break-even point rather than confining the offline cost to a discussion section.
  \item The online search cost is also misstated. Table 5 reports Stage I as 0.53 GPU-hours and the full pipeline (Stage I + Stage II with 500 candidates) as 1.01 GPU-hours, consistent with Sec. 5.4.3. However, the abstract reports 0.48 GPU-hours, and Table 4's Stage-1 column says ``0.48 for once.'' The 0.48 GPU-hours figure corresponds only to Stage II with 500 candidates.
  \item Table 3 reports search times of 0.85 and 3.84 GPU-hours for the ImageNet-1000 models D and E. These values match neither 0.48 nor 1.01 GPU-hours, and no explanation is provided for the discrepancy; Fig. 1(b) also uses 3.84 GPU-hours. Consequently, five different search-cost values (0.48, 0.53, 0.85, 1.01, and 3.84 GPU-hours) appear for the same method without reconciliation.
\end{enumerate}
Relatedly, Contribution 1 describes the search as ``training-free.'' A pipeline that trains 750 sub-networks for 25 epochs to fit a surrogate is predictor-based, not training-free, a term generally reserved in NAS for zero-cost proxies. The phrase should be removed or justified.
\response{Major Comment 1}
\begingroup
\normalfont\color{responseblue}
We thank the editor for identifying the inconsistent search-cost accounting. We agree that the original presentation conflated the platform-specific Stage-II cost with the complete online search cost and did not make the offline preparation overhead sufficiently visible. We have revised the cost presentation and the scope of the efficiency claim as follows.

\textbf{(a) Separate offline preparation and online search costs, and report the break-even point.} We have introduced a dedicated subsection, ``Results on search cost,'' immediately after the cross-platform results. Its search-cost table separately reports the offline preparation cost, the first online search, subsequent online searches, and cumulative costs for one and three platforms. For FlashNAS-BNN, the offline preparation cost of 1363.20 GPU-hours includes both generating accuracy labels for the 750 sampled architectures and fitting CIF-Net on those labeled architectures. The fitting time is short relative to label generation and is already included in the reported total; it is not an omitted additional cost. We have made this scope explicit in Section 4.2.1 and in the search-cost table, whose column is now labeled ``Offline prep.'' rather than ``Label gen.''. Stage I requires 0.53 GPU-hours and is reused across target platforms, while Stage II with 500 candidates requires 0.48 GPU-hours per platform. Thus, under the reported accounting scheme, the first online search costs 1.01 GPU-hours, and the cumulative cost over $M$ platforms is
\[
T_{\mathrm{ours}}(M)=1363.20+0.53+0.48M
                       =1363.73+0.48M\quad\text{GPU-hours}.
\]
For NAS-BNN, the corresponding reported per-platform search cost is 480 GPU-hours, giving $T_{\mathrm{NAS\text{-}BNN}}(M)=480M$. The revised table caption explicitly distinguishes this search cost from NAS-BNN's separately reported supernet-training cost of 220 GPU-days, which is excluded from the table's Total column. It also states the independent-search assumption used to extrapolate the cell-based baselines and identifies omitted, unreported offline costs. The table therefore makes its accounting scope explicit rather than presenting the comparison as fully inclusive of all training costs.

The break-even point under this scheme is
\[
M^{*}=\frac{1363.73}{480-0.48}\approx2.84.
\]
Accordingly, we now state that FlashNAS-BNN has a lower cumulative cost only when applied to at least three target platforms under this reuse setting. The table reports 1364.21 versus 480.00 GPU-hours for one platform, and 1365.17 versus 1440.00 GPU-hours for three platforms. These values acknowledge that the offline preparation overhead makes our method more expensive for a single platform. This analysis concerns cross-platform reuse for a fixed dataset and search setting; it does not assume that one CIF-Net predictor serves both ImageNet and CIFAR, for which the revised manuscript specifies separate predictors. We now explicitly state that, when a new dataset requires a separately trained predictor, its offline preparation cost must be accounted for again; the amortization argument does not assume unrestricted reuse across datasets.

\textbf{(b) Correct the meaning of 0.48 GPU-hours.} The 0.48 GPU-hours value denotes Stage II with 500 candidates, not Stage I or the full first online search. The stage-wise ablation table retains Stage I at 0.53 GPU-hours, and the CLE discussion states that Stage I plus Stage II takes 1.01 GPU-hours. We have removed the misleading ``0.48 for once'' entry and the search-time columns from the cross-platform model-comparison table, and consolidated the cost comparison in the dedicated search-cost table. The Abstract and Conclusion no longer advertise 0.48 GPU-hours as the complete search cost per platform.

\textbf{(c) Remove the inconsistent model-specific cost comparisons.} The previous 0.85 and 3.84 GPU-hours entries for the ImageNet-1000 models D and E have been removed from the model-comparison table and its accompanying discussion. The former cost-comparison panel in Fig.~1(b), which used 3.84 GPU-hours, has also been removed; the revised Fig.~1 illustrates the cell-based and chain-structured search spaces. The ImageNet results now focus on accuracy and computational budgets, while search efficiency is discussed using the explicit offline/online accounting above. We no longer use the old model-specific entries to support a search-speedup claim.

\textbf{Terminology and contribution scope.} We have removed the ``training-free'' description from Contribution 1. The revised contribution explicitly states that CIF-Net is trained using accuracy labels from sampled sub-networks. During online search, this predictor avoids candidate-by-candidate sub-network training, and the two-stage strategy reuses the accuracy-screened candidate set across platforms. The Introduction has been revised consistently to describe this predictor-based search mechanism.

Please refer to \revloc{the Abstract, Introduction, Contribution 1, and revised Fig. 1 (Page \underline{\hspace{1cm}})}, \revloc{Section 4.2.1, ``Training of accuracy predictor'' (Page \underline{\hspace{1cm}})}, \revloc{the subsections ``Results on ImageNet,'' ``Cross-platform results,'' and ``Results on search cost,'' including their comparison tables (Page \underline{\hspace{1cm}})}, and \revloc{the stage-wise ablation table, the ``CLE'' discussion, and the Conclusion (Page \underline{\hspace{1cm}})}.
\par\endgroup
\finishresponse

\item[]\textbf{Major Comment 2:} \textbf{The accuracy surrogate may be validated against the wrong target.}
CIF-Net is trained on accuracies measured after only 25 training epochs (Sec. 4.2.1), whereas the final architectures are trained using the two-phase, 256-epoch ReActNet recipe (Sec. 4.2.3), and CIFAR models are trained for 600 epochs using a different direct-quantization recipe. Table 6 reports MSE, Kendall correlation, and $R^2$, but the manuscript does not state whether these metrics are computed against the 25-epoch labels or fully trained accuracy. If they are computed against the 25-epoch labels, as implied by the 80\%/20\% split of the 750 labeled samples, then the reported Kendall correlation of 0.874 measures agreement with short-training accuracy rather than the deployment-relevant accuracy. The Stage-I screening may therefore rank architectures using a signal that does not transfer. A rank-correlation analysis between 25-epoch and fully trained accuracy on a held-out set is the minimum evidence needed; without it, the central claim of search efficiency without accuracy loss is unsupported.

Additionally, Sec. 4.2.1 does not identify the dataset or datasets used to train the 750 sampled architectures, while Table 6 reports predictor results on both ImageNet and CIFAR. It is unclear whether one CIF-Net is used for all datasets or separate predictors and labeling budgets are used.
\response{Major Comment 2}
\begingroup
\normalfont\color{responseblue}
Thank you for identifying the distinction between predicting short-training accuracy and ranking architectures by their final accuracy. A high correlation with 25-epoch labels alone does not establish that the Stage-I ranking transfers to fully trained models.

\textbf{(a) Predictor targets and data provenance.} The main manuscript's predictor metrics compare CIF-Net outputs with measured 25-epoch accuracy labels, not fully trained accuracy. We now state this explicitly alongside the predictor comparison. We train CIF-Net separately for the ImageNet and CIFAR experiments, using 25-epoch accuracy labels from ImageNet-100 and CIFAR-100, respectively. Each predictor uses labels from 750 sampled architectures and an 80\%/20\% training/held-out split; offline preparation is accounted for separately for each predictor. These clarifications appear in \revloc{Section 4.2.1 and the CIF-Net ablation discussion of the main manuscript, and Supplementary Section 3, ``CIF-Net training and cross-dataset evaluation'' (Page \underline{\hspace{1cm}})}.

\textbf{(b) Short-training versus full-training ranking.} Supplementary Table 3 reports evaluations at 15, 20, 25, and 30 epochs for three independently sampled architecture sets for each source dataset (ImageNet-100 and CIFAR-100), each containing 100 architectures and excluding the predictor's training and held-out partitions. Architecture-sampling seeds are 2024, 1234, and 42. Metrics are computed within each set and summarized as mean and standard deviation across the three sets. These are sampling repetitions, rather than three training repetitions of every fixed architecture. The model-training seed is reported separately in the training-hyperparameter table. For both Supplementary Tables 3 and 4, ImageNet-1000 reference accuracies use the two-phase ReActNet recipe with 256 epochs per phase (512 epochs in total). These analyses do not evaluate the 256-total-epoch ImageNet-1000 schedule. The budgets of 15, 20, 25, and 30 epochs are total short-training budgets per architecture. All short-training runs on ImageNet and CIFAR use a single W1A1 phase (binary weights and activations) for the entire budget. These evaluations retain the base training hyperparameters while shortening the training horizon and adjusting the learning-rate schedule to each budget; they are not intermediate checkpoints of a full-training run. At 25 epochs, ImageNet-100 and CIFAR-100 achieve within-dataset Kendall correlations of $0.79\,(0.03)$ and $0.77\,(0.03)$, Spearman correlations of $0.91\,(0.02)$ and $0.90\,(0.02)$, and Recall@10 values of $0.80\,(0.03)$ and $0.77\,(0.04)$, respectively; parentheses denote standard deviations. The complete results appear in \revloc{Supplementary Section 3.1, ``Short-training versus full-training accuracy,'' and Supplementary Table 3 (Page \underline{\hspace{1cm}})}.

\textbf{(c) Cross-dataset transfer and normalization.} The same table separately reports ImageNet-100 to ImageNet-1000 and CIFAR-100 to CIFAR-10 transfer. At 25 epochs, the corresponding Kendall correlations are $0.65\,(0.04)$ and $0.69\,(0.03)$, Spearman correlations are $0.82\,(0.03)$ and $0.85\,(0.03)$, and Recall@10 values are $0.67\,(0.04)$ and $0.70\,(0.04)$. Cross-dataset correlations and Top-10 overlap were lower than their within-dataset counterparts, indicating partial preservation of architecture ordering. The min--max-normalized MSE ($\mathrm{MSE}_{\mathrm{mm}}$) and normalized $R^2$ compare separately min--max-normalized accuracy vectors; they are not absolute target-accuracy prediction errors. Min--max normalization rescales objectives and does not establish a cross-dataset accuracy calibration. The 30-epoch setting improves all reported metrics, while increasing the labeling budget by 20\% relative to 25 epochs. We therefore describe 25 epochs as a labeling-budget versus rank-agreement trade-off, rather than claiming equivalence to full training. The measured-label comparison tests short-training signal transfer; it is separate from the main manuscript's evaluation of CIF-Net against short-training labels, and does not itself directly evaluate frozen CIF-Net outputs against final accuracy.

\textbf{(d) Search comparison under a reduced budget.} Supplementary Table 4 compares Stage-I searches guided by measured 25-epoch source accuracy and full-training target accuracy under matched OP constraints. Each run uses 64 initial candidates, 18 crossover offspring, 18 mutation offspring, and one evolutionary update. Both searches are repeated with seeds 2024, 1234, and 42. Final selection uses $w_{\mathrm{OP}}=1$, retaining accuracy--OPs Pareto filtering during evolution. The full-training target accuracies of the selected architectures differ by 0.11, 0.15, 0.12, and 0.09 percentage points on ImageNet-100, ImageNet-1000, CIFAR-100, and CIFAR-10, respectively. Differences are paired by seed before computing their mean and sample standard deviation. This reduced-budget comparison makes full-training guidance tractable while testing its effect on search outcomes. The revisions appear in \revloc{Supplementary Section 3.2, ``Accuracy--OPs search under a reduced budget,'' and Supplementary Table 4 (Page \underline{\hspace{1cm}})}.

\textbf{(e) Model-selection settings and scope.} The supplementary algorithm retains the hardware-aware accuracy--latency objective. We distinguish its weight $w_{\mathrm{Lat}}$ from an accuracy--OPs candidate-selection coefficient $w_{\mathrm{OP}}$. Reference models A--F are separated from the CPU/GPU/FPGA selections G/H/I, which use $w_{\mathrm{Lat}}=0.5$. Supplementary Table 6 explicitly reports the confirmed accuracy--OPs coefficients for A--F; F uses $w_{\mathrm{OP}}=1$ without Stage-II latency refinement. Matching F's reported performance with the CPU accuracy-only endpoint is not treated as evidence of identical search provenance. The model mapping and the CPU Stage-II weight ablation are presented in \revloc{Supplementary Sections 2.1 and 6, and Supplementary Tables 5--6 (Page \underline{\hspace{1cm}})}.
\par\endgroup
\finishresponse

\item[]\textbf{Major Comment 3:} \textbf{The cross-platform comparison is largely against manually designed models.}
In Table 4, all NAS baselines (BNAS, EBNAS, and NAS-BNN) are marked N/A on GPU and FPGA. Sec. 5.3 explains that EBNAS's group convolution lacks kernels on FPGA and GPU, but this does not explain why BNAS or NAS-BNN are unavailable: they are ReActNet-style chain architectures using standard convolutions supported by Bolt on Mali GPUs and potentially implementable on the authors' FPGA accelerator. As a result, the GPU latency reduction claim (16.0\%) and FPGA claim (6.8\%) are supported only against ReActNet-A and A\&B BNN, while the abstract's claim of consistent latency optimization across CPU, GPU, and FPGA relative to NAS methods is supported only on CPU. The manuscript also does not state whether the FPGA baseline results use the same accelerator implementation as the proposed models, which is necessary for a fair comparison.

Separately, the two NAS-BNN rows in Table 4 (430.29 ms without CLE and 272.40 ms with CLE, at 487.2 GPU-hours) are not explained. The manuscript should clarify what the second row represents, including whether NAS-BNN was rerun with latency guidance. If so, the abstract's ``up to 61.8\% over NAS-BNN'' claim uses the weaker first row (430.29 to 164.57 ms); relative to the 272.40 ms row, the reduction is 39.6\%.

Finally, the ImageNet-100 accuracies of the baselines (for example, NAS-BNN at 81.48\%) differ from the values in the original papers, which report ImageNet-1K. The manuscript does not state whether all baselines were retrained by the authors using an identical recipe.
\response{Major Comment 3}
\begingroup
\normalfont\color{responseblue}
Thank you for pointing out the need to clarify the baseline implementations, accuracy sources, and comparison scope. We have revised these descriptions as follows.

\textbf{Cross-platform baseline results.} Table 4 now includes GPU and FPGA measurements for BNAS and all four NAS-BNN variants: without CLE and with CLE targeting CPU, GPU, or FPGA. EBNAS remains unavailable on these two platforms because its group convolutions lack corresponding kernels in the implementations used. The CLE target labels identify the platform latency objective used to select each NAS-BNN architecture, and each architecture is evaluated on all three platforms. Relative to the corresponding CPU-, GPU-, and FPGA-targeted NAS-BNN variants with CLE, FlashNAS-BNN-G, H, and I reduce mean latency by 39.6\%, 10.6\%, and 8.6\%, respectively. The GPU comparison is 48.45 versus 54.20\,ms, and the FPGA comparison is 2.340 versus 2.560\,ms. These results extend the comparison beyond manually designed baselines; measured standard deviations are reported in the table. The revisions appear in \revloc{Section 5.3, ``Cross-platform results,'' and Table 4 of the main manuscript (Page \underline{\hspace{1cm}})}.


\textbf{FPGA implementation.} The FPGA baselines and FlashNAS-BNN use the same accelerator implementation. Their latency measurements follow the common protocol specified in the supplementary material. This clarification is now stated in \revloc{Section 5.3, ``Cross-platform results,'' of the main manuscript (Page \underline{\hspace{1cm}})} and \revloc{Section 1.3, ``Latency measurement protocol,'' of the supplementary material (Page \underline{\hspace{1cm}})}.

\textbf{Accuracy sources.} The ImageNet-100 accuracies of BNAS, EBNAS, and NAS-BNN are obtained by training the models on ImageNet-100 using their respective codebases and method-specific training procedures. They are not transferred from the original ImageNet-1K results. We therefore do not describe these experiments as using an identical training recipe across methods. The clarification appears in \revloc{Section 5.3 of the main manuscript (Page \underline{\hspace{1cm}})} and \revloc{Section 1.3 of the supplementary material (Page \underline{\hspace{1cm}})}.

\textbf{NAS-BNN comparison.} We distinguish NAS-BNN without CLE from the CPU-targeted variant with CLE. Their CPU latencies are $430.29\pm4.30$ and $272.40\pm2.40$\,ms, respectively. FlashNAS-BNN-G achieves $164.57\pm1.37$\,ms, corresponding to reductions of 61.8\% and 39.6\%. To avoid highlighting only the weaker baseline, the abstract and contribution summary now report the 39.6\% reduction relative to CPU-targeted NAS-BNN with CLE, at comparable accuracy. Section 5.3 reports both CPU comparisons explicitly. The revised locations are \revloc{the Abstract, contribution summary, and Section 5.3 of the main manuscript (Page \underline{\hspace{1cm}})}.
\par\endgroup
\finishresponse

\item[]\textbf{Major Comment 4:} \textbf{The FPGA latency model appears inconsistent with its motivation, and its evaluation is weak.}
Sec. 3.4.1 motivates CLE-FPGA using ``streaming and pipelined execution,'' but models end-to-end latency as the sum of per-layer latencies in Eq. 15, which is a sequential execution model. For a pipelined dataflow accelerator, steady-state throughput is governed by the slowest pipeline stage rather than the sum of stage latencies; summation is appropriate only for non-overlapped, layer-by-layer execution. Conversely, Sec. 3.4.2 justifies the CPU/GPU sequence model by stating that CPU/GPU operators execute sequentially without inter-operator pipelining, which is the regime in which a per-layer sum would be natural. The two rationales appear to be swapped, or the FPGA accelerator may in fact be non-pipelined, in which case the motivation text is inaccurate. This issue needs resolution because it determines whether Eq. 15 is a principled estimator or a regression crutch.

The empirical evidence does not resolve this concern. In Table 7, CLE has the worst within-bound prediction accuracy among the four methods on FPGA (e.g., 20.8\% versus 12.4--14.4\% at the 1\% bound), while the claimed selection advantage (2.30 ms versus 2.31/2.32 ms) is a 0.4\% latency difference reported without measurement uncertainty.
\response{Major Comment 4}
\begingroup
\normalfont\color{responseblue}
We thank the editor for identifying the inconsistency. The manuscript incorrectly interchanged the descriptions and equations of the two platform-specific estimators. In the experiments, FPGA latency was estimated using the LSTM sequence encoder with global architecture features, whereas CPU/GPU latency was estimated using the additive model. The error-bound percentages and selected-architecture latencies reported in Table 7 were obtained with this assignment. We have corrected the method descriptions, equations, and Figure 6 to match the experimental implementation; this correction does not change the underlying experiments or their reported values.

\textbf{Execution model and prediction target.} Section 3.4.1 now formulates CLE-FPGA as an LSTM layer-sequence encoder fused with global architecture features, with a single end-to-end latency output. The previous FPGA sum has been removed. The revised text distinguishes the first-result latency, the initiation interval governing steady-state throughput, and the finite-batch ms/image target. In particular, the elapsed time for a batch divided by its size includes fill/drain effects and cannot generally be replaced by either the sum of isolated full-layer times or the slowest-stage time alone. The LSTM is presented as an empirical surrogate, not as an exact analytical model of the pipeline.

\textbf{CPU/GPU formulation and relation to nn-Meter.} Section 3.4.2 now places the additive estimator under the explicit assumption of non-overlapped execution units. Units follow backend execution granularity, including fusion; the text no longer states that CPU/GPU hardware universally lacks inter-operator overlap. Figure 6 has been replaced with a schematic consistent with these two formulations. The revised discussion also recognizes nn-Meter's fusion-aware kernel prediction and its stated limitation for concurrent execution. Our FPGA formulation removes the independent-unit additivity constraint, but we do not claim that nn-Meter cannot be extended to FPGA or that the sequence model is automatically superior.

\textbf{Interpretation of the reported FPGA results.} The within-bound metric is the proportion of predictions within the specified relative-error tolerance, so larger values are better. Thus, the reported 20.8\% at the 1\% bound is higher, rather than worse, than 12.4\%--14.4\%. The current table reports CLE at 20.8\%/46.0\%/92.4\% for the 1\%/5\%/10\% bounds, and reports selected-architecture latency with standard deviation. These results evaluate the implemented platform-specific estimators. They do not by themselves establish statistical significance for small selected-latency differences or isolate the contribution of the LSTM and global-feature branches.

The formulation and scope corrections are provided in \revloc{Section 3.4, ``Cross-platform latency estimation,'' including Sections 3.4.1--3.4.2 and Figure 6 (Page \underline{\hspace{1cm}})}. Metric definitions and the reported FPGA values are provided in \revloc{Table 7 and the CLE discussion (Page \underline{\hspace{1cm}})}.
\par\endgroup
\finishresponse

\item[]\textbf{Major Comment 5:} \textbf{The latency measurement protocol is missing.}
Table 4's headline latency values (e.g., 164.57 ms on CPU and 2.34 ms on FPGA) are reported without batch size, CPU thread count, clock configuration, warm-up/repetition counts, or run-to-run variance. The FPGA improvement (2.34 versus 2.51 ms, or 6.8\%) and the Table 7 FPGA selection margins (0.01--0.02 ms) may lie within the noise of uncontrolled measurements. Given that the paper's central thesis is that measured latency, rather than proxies, should drive architecture choice, the measurement methodology must be reproducible.
\response{Major Comment 5}
\begingroup
\normalfont\color{responseblue}
We thank the editor for requesting a reproducible latency measurement protocol and explicit measurement uncertainty. We have clarified the measurement protocol and updated the repeated-search latency measurements in the supplementary material and added the measured standard deviations to the main manuscript's cross-platform comparison table. We also provide the individual-search results underlying the CLE evaluation, distinguishing timing variability from variability across searches.

\textbf{(a) Workload and measurement protocol.} Section 1.3 of the supplementary material specifies an inference batch size of 32, an input resolution of $224\times224$, 20 warm-up iterations, and 50 timed iterations per architecture. Each batch inference time is divided by 32 to obtain latency per image, and the mean and standard deviation are computed over these timed measurements. Synchronization is performed before execution and after inference completes. The timing boundary covers model inference from a prepared input tensor to the model output, excluding host-side image preprocessing. This protocol applies to all measured inference latencies in the main manuscript and supplementary material. The main manuscript also distinguishes the batch size of 1 used for latency-predictor training from the inference batch size.

\textbf{(b) Software, execution settings, and cooling.} CPU and GPU inference use Huawei's Noah's Bolt v1.5.2. The ARM Cortex-X4 CPU uses a single inference thread with Bolt's \texttt{CPU\_AFFINITY\_HIGH\_PERFORMANCE} policy, and the ARM Mali-G77 MP9 GPU uses Bolt's default OpenCL backend. The CPU and GPU are configured to their highest frequency operating points of 3.3~GHz and 836~MHz, respectively. The FPGA implementation targets the Xilinx XC7VX690T, with its clock configured to 180~MHz, and uses Vitis HLS 2022.1 and Vivado 2022.1. Its implementation basis is now identified through the reference to Zhang et al.~\cite{Zhang2025}. An external fan is used during benchmarking to reduce thermal effects on the measurements.

\textbf{(c) Measurement uncertainty and repeated-search results.} The cross-platform comparison table now reports measured latency as mean $\pm$ standard deviation. For example, the CPU latency of FlashNAS-BNN-G is $164.57\pm1.37$~ms. On FPGA, FlashNAS-BNN-I achieves $2.340\pm0.008$~ms, compared with $2.510\pm0.008$~ms for ReActNet-A. The mean difference of 0.170~ms (6.8\%) is larger than the reported timing standard deviation of either model.

For the CLE evaluation, Section 7.2 of the supplementary material reports three searches per method and platform, using seeds 2024, 1234, and 42 for Search 1, 2, and 3, respectively. The Actually measured baseline uses actual hardware latency to guide architecture search and follows the same repeated-search setting. Each Search column reports the selected architecture's measured mean $\pm$ timing standard deviation. The Overall column instead reports the mean $\pm$ sample standard deviation of the three search-specific mean latencies; these Overall values are also reported in the main manuscript. The error-bound percentages are summarized as mean $\pm$ sample standard deviation over the same three searches.

CLE yields Overall latencies of $141.32\pm0.53$~ms on CPU, $43.91\pm0.12$~ms on GPU, and $2.300\pm0.005$~ms on FPGA. The corresponding Actually measured results are $139.97\pm0.42$~ms, $43.78\pm0.04$~ms, and $2.297\pm0.004$~ms. The FPGA mean latency gap is 0.003~ms. These results show that CLE remains close to search guided by actual hardware latency. Across all seven FPGA methods in Supplementary Table 7, the per-search timing SDs range from 0.005 to 0.012~ms and the between-search sample SDs range from 0.004 to 0.075~ms. These describe different sources of variability and are not confidence intervals or tests of significance. We therefore do not interpret the small FPGA mean-latency differences as statistically significant.

The revisions are provided in \revloc{Section 4.2.2, ``Training of latency predictor,'' the cross-platform comparison table, and the CLE evaluation table and discussion in the main manuscript (Page \underline{\hspace{1cm}})}, \revloc{Section 1.3, ``Latency measurement protocol,'' in the supplementary material (Page \underline{\hspace{1cm}})}, and \revloc{Section 7.2, ``CLE performance over repeated searches,'' in the supplementary material (Page \underline{\hspace{1cm}})}.

\par\endgroup
\finishresponse

\item[]\textbf{Major Comment 6:} \textbf{The scope of the claims and internal consistency of the results need attention.}
\begin{enumerate}[label=(\alph*),leftmargin=2em,itemsep=2mm]
  \item The abstract claims latency reductions ``up to 47.1\% over ReActNet-A and up to 61.8\% over NAS-BNN, without accuracy loss,'' whereas the conclusion states ``6.1\%--47.1\%.'' The 61.8\% figure is absent from the conclusion, and 6.1\% does not match a value in Table 4; the FPGA reduction is 6.8\%.
  \item The phrase ``without accuracy loss'' is not literally supported for the headline CPU model: FlashNAS-BNN-G achieves 81.40\%, below NAS-BNN at 81.48\% in the same table.
  \item Fig. 7(b) plots ``FlashNAS-BNN-I 1x1 Conv'' points, while Sec. 5.3 describes FlashNAS-BNN-I as an ``all-3x3 convolution design.'' Residual pointwise convolutions may explain the difference, but this should be stated.
  \item Sec. 5.4.2 claims a 24x speedup over standalone EA when CIF-Net is combined with two-stage search. Table 5 gives 22.81/0.53 = 43x for that comparison, while 24x corresponds to a different pair of values (550.53/22.81).
  \item Sec. 2.1 lists four search times (96, 9.6, 6, and 0.96 GPU-hours) for three methods. Although each discrepancy may be minor, together they suggest that the numbers have not been cross-checked, which is consequential in a paper whose contribution is quantitative.
\end{enumerate}
\response{Major Comment 6}
\begingroup
\normalfont\color{responseblue}
Thank you for identifying these inconsistencies. We have checked the comparison baselines, accuracy claims, speedup calculations, and reported search times. We address the individual points below.

\begin{enumerate}[label=(\alph*),leftmargin=2em,itemsep=2mm]
\item \textbf{Latency reductions and comparison baselines.} The abstract now reports the reductions separately for CPU, GPU, and FPGA: 47.1\%/16.0\%/6.8\% relative to ReActNet-A and 39.6\%/10.6\%/8.6\% relative to the corresponding platform-optimized NAS-BNN architectures with CLE. The previous 61.8\% CPU reduction used NAS-BNN without CLE as the baseline; Section 5.3 now distinguishes this comparison from the 39.6\% reduction relative to CPU-targeted NAS-BNN with CLE. We have also corrected the conclusion to report 6.8\%--47.1\% relative to ReActNet-A and 8.6\%--39.6\% relative to platform-optimized NAS-BNN, explicitly identifying both comparison groups. These changes appear in \revloc{the Abstract, Section 5.3, ``Cross-platform results,'' and the Conclusion (Page \underline{\hspace{1cm}})}.

\item \textbf{Accuracy wording.} We have replaced ``without accuracy loss'' with ``at comparable accuracy.'' FlashNAS-BNN-G achieves 81.40\% Top-1 accuracy, compared with 81.48\% for NAS-BNN without CLE, a difference of 0.08 percentage points. We therefore do not claim that accuracy is unchanged. The revised wording accompanies the latency comparisons in \revloc{the Abstract, the contribution summary, Section 5.3, and the Conclusion (Page \underline{\hspace{1cm}})}.

\item \textbf{Architecture description and figure legend.} The searched FlashNAS-BNN-I architecture uses an all-$3\times3$ convolution design as described in Section 5.3. The ``FlashNAS-BNN-I 1x1 Conv'' legend entry was retained in error and had no corresponding data points in Fig.~7(b). Thus, this discrepancy concerns the figure legend rather than residual pointwise convolutions in the searched architecture. We have removed the erroneous legend entry and replaced Fig.~7(b) in the revised manuscript with the corrected plot. The ReActNet $1\times1$ and $3\times3$ convolution entries and the FlashNAS-BNN-I $3\times3$ convolution entry are retained. The correction is provided in \revloc{Figure 7(b), Section 5.3, ``Cross-platform results,'' in the main manuscript (Page \underline{\hspace{1cm}})}.

\item \textbf{Speedup calculations.} Section 5.5.2 now identifies the two comparisons separately: CIF-Net alone reduces the standalone EA Stage-I cost from 550.53 to 22.81 GPU-hours (approximately $24\times$), and adding the two-stage strategy further reduces it from 22.81 to 0.53 GPU-hours (approximately $43\times$). These are online Stage-I search costs; the offline preparation cost is separately reported in ``Results on search cost. The corrected comparisons are provided in \revloc{Section 5.5.2, ``Two-stage Search,'' and Table 5 (Page \underline{\hspace{1cm}})}.

\item \textbf{Search-time correspondence.} Section 2.1 now lists three search times for the three methods: 9.6 GPU-hours for BNAS, 6 GPU-hours for BATS, and 0.96 GPU-hours for EBNAS. The extraneous 96 GPU-hours entry has been removed. The correction is provided in \revloc{Section 2.1, ``Neural architecture search for BNNs'' (Page \underline{\hspace{1cm}})}.
\end{enumerate}
\par\endgroup
\finishresponse

\item[]\textbf{Major Comment 7:} \textbf{Baselines and methodological details needed to assess attribution are missing.}
\begin{enumerate}[label=(\alph*),leftmargin=2em,itemsep=2mm]
  \item The motivating claim is that OPs are a poor proxy for latency (Fig. 2), yet Table 7 does not include the simple baseline of ranking the same candidate set by OPs. If OPs ranking also selects a near-optimal architecture on these platforms, CLE's practical value would be substantially reduced.
  \item The $(w_1,w_2)$ settings used to produce models D--I are not given in the main text and are deferred to supplementary material 1.5, which was unavailable for review. The operating points of the headline models therefore cannot be interpreted.
  \item nn-Meter, a canonical work on per-platform latency prediction, is neither cited nor compared. The compared latency predictors, BRP-NAS and CLCDLP, are cell-based GNN methods that the manuscript itself argues are ill-suited to chain-structured search spaces, raising concern about the choice of baselines.
  \item The two-stage search algorithm is only summarized in supplementary material 1.1. The main text gives the objective (Eqs. 4--8) but does not describe the evolutionary procedure in sufficient detail for reproduction.
\end{enumerate}
\response{Major Comment 7}
\begingroup
\normalfont\color{responseblue}
Thank you for identifying the missing comparisons and reproducibility details. We have added the OPs Proxy and nn-Meter results, separated the two selection weights, and expanded the algorithm and configuration descriptions.

\textbf{(a) OPs baseline and comparison scope.} The main manuscript's CLE evaluation table now includes OPs Proxy on all three platforms. Its selected-architecture Overall latencies are $151.78\pm5.73$, $46.07\pm2.19$, and $2.394\pm0.075$~ms/image on CPU, GPU, and FPGA, respectively, compared with $141.32\pm0.53$, $43.91\pm0.12$, and $2.300\pm0.005$~ms/image for CLE. These are mean $\pm$ sample SD across three searches. Supplementary Table 7 provides the search-specific measurements. The estimator comparison uses matched OP constraints and latency-only Stage-II selection, $w_{\mathrm{Lat}}=0$, separately from the default accuracy--latency selection for G/H/I. For each seed, OPs Proxy and CLE reuse exactly the same candidate architectures. OPs Proxy ranks this shared pool by operation count, while CLE uses estimated latency; the selected architectures are evaluated using measured hardware latency.

\textbf{(b) Weights and model mapping.} We now use one accuracy weight per objective: $w_{\mathrm{OP}}$ for accuracy--OPs selection and $w_{\mathrm{Lat}}$ for Stage-II accuracy--latency selection; the other objective receives weight $1-w$. Stage-I non-dominated filtering itself has no scalar weight. Supplementary Table 6 reports $w_{\mathrm{OP}}=0.25,0.50,0.75$ for A$^*$, B$^*$, C$^*$, and $0.50,0.75,1.00$ for D, E, F, respectively. These reference variants do not use Stage-II latency refinement. G/H/I use $w_{\mathrm{Lat}}=0.5$ for CPU/GPU/FPGA, respectively, with platform-specific latency normalization. N/A denotes an objective that was not applied. Supplementary Table 5 reports the CPU weight sweep. Its $w_{\mathrm{Lat}}=0.5$ result matches G's 81.40\% accuracy and $164.57\pm1.37$~ms/image CPU latency; the accuracy-only endpoint matches F's reported 82.07\% and $412.91\pm3.80$~ms/image, without establishing identical search provenance. For all five weight settings, measured CPU latency is reported as mean $\pm$ SD over 50 timed iterations. The results for weights 0, 0.25, and 0.75 are $139.97\pm1.15$, $164.57\pm1.37$, and $255.35\pm2.12$~ms/image, respectively. The timing SD of 1.15~ms at weight 0 is distinct from the between-search SD of 0.42~ms in the Actually measured comparison.

\textbf{(c) nn-Meter and graph-based baselines.} We added nn-Meter to the related-work discussion, bibliography, and three-platform estimator comparison. Its Overall latencies are $142.86\pm0.81$, $44.17\pm0.19$, and $2.309\pm0.007$~ms/image on CPU, GPU, and FPGA, respectively. The revised modeling discussion distinguishes additive CPU/GPU execution from overlapping FPGA pipeline stages. The expanded comparison includes both nn-Meter and the simple OPs Proxy baseline alongside BRP-NAS and CLCDLP, allowing CLE to be assessed against alternatives beyond graph-based predictors.

\textbf{(d) Evolutionary procedure and supplementary access.} Main-text Section 3.2 now describes candidate representation, initialization, crossover and mutation, CIF-Net evaluation, non-dominated filtering, and termination. Section 4.2.3 reports the numerical search settings and model-specific weights. Supplementary Section 1.2 additionally lists the mutation probability and candidate selection size. Algorithm 1 in Supplementary Section 2 details the two-stage workflow, including the direct-measurement/CLE branch; Supplementary Section 7.1 explains the threshold of 250 candidates. The complete supplementary document is also appended to this response.

The revisions are located in \revloc{Section 2 (related work), Sections 3.2 and 3.4, Section 4.2.3 (search settings and model-specific weights), and the CLE evaluation discussion and table in the main manuscript; Supplementary Sections 1.2, 2, 6, 7.1, and 7.2, and Tables 5--7 (Page \underline{\hspace{1cm}})}.
\par\endgroup
\finishresponse
\end{enumerate}
\begin{enumerate}[leftmargin=0pt,label={},itemsep=5mm]
\bfseries
  \item[]\textbf{Minor Comment 1:} Sec. 5.4 states that the ablation ``quantitatively reveals individual contributions of these components to enhancing the detector's overall capability.'' This is a classification paper; the sentence appears to have been carried over from a detection manuscript and should be rewritten.
  \response{Minor Comment 1}
\begingroup
\normalfont\color{responseblue}
Thank you for identifying this incorrect task description. We have replaced the reference to the detector's overall capability with: ``This ablation study quantifies the contribution of each component to reducing architecture search time for image classification.'' The revised sentence describes both the classification task and the search-time quantities evaluated in the ablation study. The change appears in \revloc{the opening paragraph of ``Ablation Study'' in the main manuscript (Page \underline{\hspace{1cm}})}.
\par\endgroup
\finishresponse
  \item[]\textbf{Minor Comment 2:} Reference [12] (Krizhevsky, 2009, ``Learning multiple layers of features from tiny images'') is cited for ImageNet-1000 in Sec. 4.1, but it is the CIFAR reference. ImageNet should cite Deng et al. (2009). References [12] and [13] are currently duplicates.
  \response{Minor Comment 2}
\begingroup
\normalfont\color{responseblue}
Thank you for pointing out the incorrect citation and duplicate references. Section 4.1 now cites Deng et al. (2009), ``ImageNet: A Large-Scale Hierarchical Image Database,'' for ImageNet-1000, while Krizhevsky (2009), ``Learning multiple layers of features from tiny images,'' is retained for CIFAR-10/100. We have removed the duplicate CIFAR bibliography entry and consolidated the citations under a single reference. These corrections are provided in \revloc{Section 4.1, ``Dataset,'' and the References of the main manuscript (Page \underline{\hspace{1cm}})}.
\par\endgroup
\finishresponse
  \item[]\textbf{Minor Comment 3:} The CIFAR improvements over EBNAS in Table 2 (94.16 versus 94.04 and 73.10 versus 72.99) are within typical run-to-run variance for CIFAR BNN training. Without repeated runs or variance estimates, ``matching or slightly exceeding'' is the strongest defensible claim. The search-time advantage over EBNAS is also nil (0.96 versus 0.96 GPU-hours).
  \response{Minor Comment 3}
\begingroup
\normalfont\color{responseblue}
We agree that the small CIFAR accuracy differences do not establish a statistically significant improvement. The discussion now describes FlashNAS-BNN-C as ``matching or slightly exceeding'' EBNAS and explicitly states that the differences of 0.12 and 0.11 percentage points cannot establish significance without repeated-run statistics. We also state that FlashNAS-BNN does not demonstrate a search-time advantage over EBNAS for a single search. The cost comparison is presented separately in the revised search-cost table, with offline preparation and online search distinguished as explained in our response to Major Comment 1. The revisions are provided in \revloc{the discussion following the CIFAR comparison table in Section 5.1, ``Results on CIFAR,'' and the ``Results on search cost'' subsection of the main manuscript (Page \underline{\hspace{1cm}})}.
\par\endgroup
\finishresponse
  \item[]\textbf{Minor Comment 4:} Table 3 contains the typographical errors ``Mannal'' and ``Train epoches.'' Table 4 inconsistently uses ``N/A'' and ``NA.''
  \response{Minor Comment 4}
\begingroup
\normalfont\color{responseblue}
Thank you for identifying these typographical inconsistencies. We have corrected ``Mannal'' to ``Manual,'' standardized the ImageNet comparison table heading as ``Training epochs,'' and unified unavailable latency entries as ``N/A.'' The corrections appear in \revloc{the CIFAR and ImageNet comparison tables and the cross-platform comparison table in the main manuscript (Page \underline{\hspace{1cm}})}.
\par\endgroup
\finishresponse
  \item[]\textbf{Minor Comment 5:} The roofline derivation in Eq. 2 assumes $C_{in}=C_{out}$ and activation dominance, $HW\gg C_iK_i^2$. This approximation may be reasonable for early stages, but it is used to motivate kernel and width choices network-wide. In addition, Fig. 4(a) validates the argument only on FPGA, while the argument is invoked for all platforms.
  \response{Minor Comment 5}
\begingroup
\normalfont\color{responseblue}
We agree that the original simplification required a limited interpretation. We have replaced the equal-channel expression with a general-channel structural intensity proxy and explicitly stated the equal-spatial-size and equal-storage-cost assumptions. The activation-dominance and equal-channel simplifications are summarized in the main text and derived in the supplementary material. The main manuscript now distinguishes this proxy from device-level Ops/Byte and states that the device Roofline bound requires consistent operation and byte units. The supplementary section ``Structural intensity assumptions'' provides the byte-aware expression for unequal input/output spatial dimensions and explains the limits of ideal-reuse traffic estimates for BNN implementations.

We have revised the search-space motivation to distinguish the roles of the structural expression, the FPGA example, and platform-specific selection: the expression motivates searching kernel size and width through their computation--data-movement trade-off, Figure 4(a) illustrates that motivation on FPGA, and CLE with measured latency guides the actual hardware-specific choice. Higher structural intensity is not claimed to guarantee lower latency. We have also removed the unconditional description of stage depth as a linear multiplier of end-to-end latency. Figure 4(a) and its caption are explicitly FPGA-specific; the unsupported comparison that FPGA latency profiling is ``less noisy'' has been removed. Platform-specific selection remains guided by CLE and measured hardware latency, rather than by directly applying the structural proxy network-wide.

These changes clarify the analysis and its scope; they do not add new CPU/GPU Roofline measurements or change the search space and NAS results. The earlier supplementary section ``Additional motivation for CLE'' remains removed; the new assumptions section does not reinstate it as experimental evidence. The revisions appear in \revloc{Section 3.1, ``Architecture search space,'' and Figure 4 in the main manuscript, and ``Structural intensity assumptions'' in the supplementary material (Page \underline{\hspace{1cm}})}. The platform-specific estimator rationale is provided in \revloc{Section 3.4, ``Cross-platform latency estimation,'' in the main manuscript (Page \underline{\hspace{1cm}})}.
\par\endgroup
\finishresponse
\end{enumerate}

\bibliographystyle{plain}
\bibliography{Paper1}

% Include the independently compiled supplement to preserve its original
% document class, pagination, numbering, bibliography, and visual style.
\clearpage
\includepdf[pages=-,pagecommand={\thispagestyle{empty}},fitpaper=true,noautoscale=true]{supplementary/elsarticle-template-num.pdf}

\end{document}
